\section{Một trung bình là một trung bình}
\label{sec:ch07-nhieu-dau}

\index{multi-head!vì sao một head là không đủ}%
\eqref{eq:ch07-attention} trả về một tổng có trọng số. Một cái.

Đó là hạn chế đọc ra được từ chính công thức chứ không cần thí nghiệm nào. Khi
từ \code{mèo} trong một câu muốn hỏi hai chuyện khác nhau, tính từ nào thuộc về
tôi và động từ nào lấy tôi làm chủ ngữ, nó chỉ có một phân phối \(\alpha\) để
diễn đạt cả hai. Trộn hai câu hỏi vào một phân phối thì được một câu trả lời
trung bình cho một câu hỏi không ai đặt.

Bài giải bằng cách chạy song song \(h\) bản, mỗi bản hẹp hơn:

\begin{equation}
  \multihead(\mat{Q}, \mat{K}, \mat{V})
  = \left[\headop_1; \dots; \headop_h\right] \mat{W}^O,
  \qquad
  \headop_i = \attention(\mat{Q}\mat{W}^Q_i, \mat{K}\mat{W}^K_i,
                         \mat{V}\mat{W}^V_i).
  \label{eq:ch07-multihead}
\end{equation}

Mỗi head chiếu query, key và value xuống một không gian con riêng, chấm điểm ở
đó, rồi cả \(h\) kết quả nối lại và đi qua một phép chiếu cuối. Bài chạy
\(h = 8\) với \(\dk = \dv = \dmodel/h = 64\).

Chữ head cũng giữ nguyên tiếng Anh, và vì đúng lý do đã giữ nguyên query, key và
value ở mục~\ref{sec:ch07-tu-nhin}: bản dịch sát là \enquote{đầu}, mà \emph{đầu}
tiếng Việt còn nghĩa là chỗ bắt đầu, nên \enquote{tám đầu} với \enquote{đầu
tiên} đứng gần nhau đọc thành hai thứ khác hẳn.

\subsection*{Chỗ dễ đọc sai: nhiều head không tốn thêm}

\index{multi-head!số tham số không đổi theo h}%
Hình vẽ trong bài xếp \(h\) khối cạnh nhau, nên phản xạ đầu tiên là tám head
tốn gấp tám lần. Không phải. Câu của bài:

\begin{quote}
\enquote{Due to the reduced dimension of each head, the total computational
cost is similar to that of single-head attention with full dimensionality.}
\end{quote}

Lý do nằm gọn trong \(\dk = \dmodel/h\). Tăng \(h\) lên gấp đôi thì mỗi head hẹp
đi một nửa, nên tổng bề rộng của \(h\) phép chiếu vẫn đúng bằng \(\dmodel\).
Bốn ma trận \(\mat{W}^Q\), \(\mat{W}^K\), \(\mat{W}^V\), \(\mat{W}^O\) giữ
nguyên hình dạng \(\dmodel \times \dmodel\) với mọi \(h\), và số tham số không
đổi. Repo lưu chúng ở dạng gộp và tách ra từng head bằng một phép \code{view}, vì đó
đúng là \(h\) ma trận của bài xếp cạnh nhau:

\begin{minted}{python}
def _split(self, x: torch.Tensor) -> torch.Tensor:
    batch, steps, _ = x.shape
    return x.view(batch, steps, self.n_heads, self.d_k).transpose(1, 2)
\end{minted}

\code{tests/test\_ch07.py} khẳng định chuyện ấy bằng một test dựng cùng một tầng
ở \(h\) bằng 1, 2, 4, 8, 16 rồi đếm tham số, và đòi năm con số bằng nhau. Tôi
viết test ấy không phải vì nghi mã, mà vì đây là chỗ một lần sửa \enquote{cho
gọn} sẽ phá lặng lẽ: đổi \(\dk\) thành một hằng số thay vì \(\dmodel/h\) thì mọi
thứ vẫn chạy, vẫn học được, và hàng (A) của bảng 3 dưới đây mất nghĩa.

\index{multi-head!FLOP và đồng hồ không cùng nói một chuyện}%
Một cảnh báo về chữ \enquote{cost} trong câu trích. Nó đúng theo phép đếm số
phép tính dấu chấm động, floating point operations, viết tắt là FLOP,
và đó là thứ bài đang nói. Theo đồng hồ thì không hẳn: một phiên đo riêng ở
\(\dmodel = 512\) trên CPU thấy \(h = 16\) chậm hơn \(h = 1\) khoảng hai mươi
phần trăm, vì chia thành nhiều phép nhân ma trận nhỏ hơn thì phần chi phí không
phải số học lớn lên. Con số ấy phụ thuộc cài đặt và phần cứng đủ nhiều để tôi
không in nó thành một phát biểu chung, và
chương~\ref{ch:chi-phi-va-song-song} là chỗ đo thời gian chạy cho tử tế.

\subsection*{Bài tự đo số head, và kết quả không đơn điệu}

\index{multi-head!hàng (A) của bảng 3}%
Hàng (A) của bảng 3 đổi \(h\) và giữ nguyên tổng chi phí tính, đúng theo cách
vừa nói. Cột \code{PPL} là perplexity và cột \code{BLEU} là điểm Bilingual
Evaluation Understudy, viết tắt là BLEU, cả hai đo trên tập phát triển
Anh-Đức:

\begin{minted}{text}
h    d_k   d_v   PPL    BLEU
1    512   512   5.29   24.9
4    128   128   5.00   25.5
16   32    32    4.91   25.8
32   16    16    5.01   25.4
\end{minted}

Nhận xét của bài, nguyên văn: \enquote{While single-head attention is 0.9 BLEU
worse than the best setting, quality also drops off with too many heads}. Hiệu
25.8 trừ 24.9 đúng bằng 0.9.

Bảng ấy nói hai chuyện và chỉ chuyện thứ nhất hay được trích. Chuyện thứ nhất:
một head là kém, đúng như lập luận ở đầu mục. Chuyện thứ hai: 32 head cũng kém,
và kém hơn cả 4 head. Ở \(h = 32\) thì \(\dk = 16\), tức mỗi head chỉ còn mười sáu
chiều để dựng câu hỏi và câu trả lời trong đó. Cơ chế đổi từ một trung bình quá
thô sang nhiều trung bình quá hẹp, và đâu đó ở giữa là chỗ tốt nhất, mà bài tìm
bằng cách thử chứ không bằng cách dẫn xuất.

Tôi thích hàng này vì nó là một bảng bốn dòng bác bỏ được cách kể \enquote{càng
nhiều head càng tốt}, và bài in nó ra ngay cạnh con số họ chọn.
